\section{Qualitative evaluation and examples}
\label{app:qualexamples}

\paragraph{Example A: \texttt{jfs\_dmap.c} boundary check.}
Listing~\ref{lst:jfs_semantic}\footnote{Example A in Syzkaller: \url{https://syzkaller.appspot.com/bug?id=5eb8a5d29d77f8a364cf3270bf9625eb4d4ffc52}}   compares the developer’s ground-truth patch  with
the patch generated by \cresearcher{}.  Both fixes add a
lower-bound check on \verb|bmp->db_agl2size| alongside the existing upper-bound
check; the only difference is the ordering of the two disjuncts in the \texttt{if}
condition, an immaterial variation in this case.
This illustrates the class of \emph{Accurate} patches.


\begin{code}
\captionof{listing}{Semantically equivalent patch
produced by \cresearcher{} for the \texttt{jfs\_dmap.c} crash.}
\label{lst:jfs_semantic}
\begin{minted}[linenos,breaklines,xleftmargin=1.5em,numbersep=5pt]{diff}
--- a/fs/jfs/jfs_dmap.c                         /* developer */
+++ b/fs/jfs/jfs_dmap.c
@@ -193,7 +193,8 @@ int dbMount(struct inode *ipbmap)
    bmp->db_agwidth = le32_to_cpu(dbmp_le->dn_agwidth);
    bmp->db_agstart = le32_to_cpu(dbmp_le->dn_agstart);
    bmp->db_agl2size = le32_to_cpu(dbmp_le->dn_agl2size);
-   if (bmp->db_agl2size > L2MAXL2SIZE - L2MAXAG) {
+   if (bmp->db_agl2size > L2MAXL2SIZE - L2MAXAG ||
+       bmp->db_agl2size < 0) {
        err = -EINVAL;
        goto err_release_metapage;
    }

--- a/fs/jfs/jfs_dmap.c                         /* generated */
+++ b/fs/jfs/jfs_dmap.c
@@ -193,7 +193,7 @@ int dbMount(struct inode *ipbmap)
    bmp->db_agwidth = le32_to_cpu(dbmp_le->dn_agwidth);
    bmp->db_agstart = le32_to_cpu(dbmp_le->dn_agstart);
    bmp->db_agl2size = le32_to_cpu(dbmp_le->dn_agl2size);
-   if (bmp->db_agl2size > L2MAXL2SIZE - L2MAXAG) {
+   if (bmp->db_agl2size < 0 ||  bmp->db_agl2size > L2MAXL2SIZE - L2MAXAG) {
        err = -EINVAL;
        goto err_release_metapage;
    }
\end{minted}
\end{code}

\paragraph{Example B: \texttt{hci\_h5.c} null-check addition.}
In Listing~\ref{lst:h5_overspec}\footnote{Example B in Syzkaller: \url{https://syzkaller.appspot.com/bug?id=1bc76a5779300075e3466a0e2abcd753df212779}}, both the developer and \cresearcher{} address an unsafe access to \texttt{hu->serdev->dev}, but the generated patch goes beyond the developer's fix. While the developer simply guards the power management calls with a null check, \cresearcher{} adds an \texttt{else} branch that logs an error and returns \texttt{-ENODEV}. This reflects a conservative design that prevents execution in the event of a null pointer, even though the surrounding kernel code may already guarantee that \texttt{hu->serdev} is non-null. Such overspecialization can be seen as a benign deviation: the patch is functionally correct and improves robustness, but at the risk of silently diverging from upstream assumptions. This illustrates the class of \emph{overspecialized} patches.

\begin{code}
\captionof{listing}{Overspecialized patch generated for \texttt{hci\_h5.c}.}
\label{lst:h5_overspec}
\begin{minted}[linenos,breaklines,xleftmargin=1.5em,numbersep=5pt]{diff}
--- a/drivers/bluetooth/hci_h5.c                /* developer */
+++ b/drivers/bluetooth/hci_h5.c
@@ -587,9 +587,11 @@ static int h5_recv(...)
    count -= processed;
    }

-   pm_runtime_get(&hu->serdev->dev);
-   pm_runtime_mark_last_busy(&hu->serdev->dev);
-   pm_runtime_put_autosuspend(&hu->serdev->dev);
+   if (hu->serdev) {
+       pm_runtime_get(&hu->serdev->dev);
+       pm_runtime_mark_last_busy(&hu->serdev->dev);
+       pm_runtime_put_autosuspend(&hu->serdev->dev);
+   }

    return 0;
 }


--- a/drivers/bluetooth/hci_h5.c                /* generated */
+++ b/drivers/bluetooth/hci_h5.c
@@ -587,9 +587,14 @@ static int h5_recv(...)
 		count -= processed;
 	}
    
-   pm_runtime_get(&hu->serdev->dev);
-   pm_runtime_mark_last_busy(&hu->serdev->dev);
-   pm_runtime_put_autosuspend(&hu->serdev->dev);
+   if (hu->serdev) {
+       pm_runtime_get(&hu->serdev->dev);
+       pm_runtime_mark_last_busy(&hu->serdev->dev);
+       pm_runtime_put_autosuspend(&hu->serdev->dev);
+   } else {
+       bt_dev_err(hu->hdev, "serdev is not initialized");
+       return -ENODEV;
+   }

    return 0;
 }
\end{minted}
\end{code}

\paragraph{Example C: \texttt{ns.c} RCU read lock insertion.}
In Listing~\ref{lst:ns_rcu_partial}\footnote{Example C in Syzkaller: \url{https://syzkaller.appspot.com/bug?id=07c9d71dc1a215b19c6a245c68f502bc57dbdb83}} both the developer and \cresearcher{} address the unsafe traversal of a radix tree without proper RCU synchronization. The developer applies a comprehensive fix, wrapping all relevant \verb|radix_tree_for_each_slot| iterations with \verb|rcu_read_lock()| and \verb|rcu_read_unlock()| across multiple functions. In contrast, \cresearcher{} focuses only on the \verb|ctrl_cmd_new_lookup()| function, inserting the necessary locking primitives in that scope alone. While this partial patch is not directly mergeable due to its incompleteness, it demonstrates an accurate understanding of the underlying concurrency issue and correctly applies the mitigation in the context it modifies. As such, it exemplifies the class of \emph{incomplete} patches, offering concrete insight into the nature and location of the bug, and accelerating the path toward a complete and upstreamable fix.

\begin{code}
\captionof{listing}{Developer and plausible patches for \texttt{ns.c}.}
\label{lst:ns_rcu_partial}
\begin{minted}[linenos,breaklines,xleftmargin=1.5em,numbersep=5pt]{diff}
--- a/net/qrtr/ns.c                         /* developer */
+++ b/net/qrtr/ns.c
@@ -193,12 +193,13 @@ static int announce_servers(struct sockaddr_qrtr *sq)
 	struct qrtr_server *srv;
 	struct qrtr_node *node;
 	void __rcu **slot;
-	int ret;
+	int ret = 0;
 
 	node = node_get(qrtr_ns.local_node);
 	if (!node)
 		return 0;
 
+	rcu_read_lock();
 	/* Announce the list of servers registered in this node */
 	radix_tree_for_each_slot(slot, &node->servers, &iter, 0) {
 		srv = radix_tree_deref_slot(slot);
@@ -206,11 +207,14 @@ static int announce_servers(struct sockaddr_qrtr *sq)
 		ret = service_announce_new(sq, srv);
 		if (ret < 0) {
 			pr_err("failed to announce new service\n");
-			return ret;
+			goto err_out;
 		}
 	}
 
-	return 0;
+err_out:
+	rcu_read_unlock();
+
+	return ret;
 }
 
 static struct qrtr_server *server_add(unsigned int service,
@@ -335,7 +339,7 @@ static int ctrl_cmd_bye(struct sockaddr_qrtr *from)
 	struct qrtr_node *node;
 	void __rcu **slot;
 	struct kvec iv;
-	int ret;
+	int ret = 0;
 
 	iv.iov_base = &pkt;
 	iv.iov_len = sizeof(pkt);
@@ -344,11 +348,13 @@ static int ctrl_cmd_bye(struct sockaddr_qrtr *from)
 	if (!node)
 		return 0;
 
+	rcu_read_lock();
 	/* Advertise removal of this client to all servers of remote node */
 	radix_tree_for_each_slot(slot, &node->servers, &iter, 0) {
 		srv = radix_tree_deref_slot(slot);
 		server_del(node, srv->port);
 	}
+	rcu_read_unlock();
 
 	/* Advertise the removal of this client to all local servers */
 	local_node = node_get(qrtr_ns.local_node);
@@ -359,6 +365,7 @@ static int ctrl_cmd_bye(struct sockaddr_qrtr *from)
 	pkt.cmd = cpu_to_le32(QRTR_TYPE_BYE);
 	pkt.client.node = cpu_to_le32(from->sq_node);
 
+	rcu_read_lock();
 	radix_tree_for_each_slot(slot, &local_node->servers, &iter, 0) {
 		srv = radix_tree_deref_slot(slot);
 
@@ -372,11 +379,14 @@ static int ctrl_cmd_bye(struct sockaddr_qrtr *from)
 		ret = kernel_sendmsg(qrtr_ns.sock, &msg, &iv, 1, sizeof(pkt));
 		if (ret < 0) {
 			pr_err("failed to send bye cmd\n");
-			return ret;
+			goto err_out;
 		}
 	}
 
-	return 0;
+err_out:
+	rcu_read_unlock();
+
+	return ret;
 }
 
 static int ctrl_cmd_del_client(struct sockaddr_qrtr *from,
@@ -394,7 +404,7 @@ static int ctrl_cmd_del_client(struct sockaddr_qrtr *from,
 	struct list_head *li;
 	void __rcu **slot;
 	struct kvec iv;
-	int ret;
+	int ret = 0;
 
 	iv.iov_base = &pkt;
 	iv.iov_len = sizeof(pkt);
@@ -434,6 +444,7 @@ static int ctrl_cmd_del_client(struct sockaddr_qrtr *from,
 	pkt.client.node = cpu_to_le32(node_id);
 	pkt.client.port = cpu_to_le32(port);
 
+	rcu_read_lock();
 	radix_tree_for_each_slot(slot, &local_node->servers, &iter, 0) {
 		srv = radix_tree_deref_slot(slot);
 
@@ -447,11 +458,14 @@ static int ctrl_cmd_del_client(struct sockaddr_qrtr *from,
 		ret = kernel_sendmsg(qrtr_ns.sock, &msg, &iv, 1, sizeof(pkt));
 		if (ret < 0) {
 			pr_err("failed to send del client cmd\n");
-			return ret;
+			goto err_out;
 		}
 	}
 
-	return 0;
+err_out:
+	rcu_read_unlock();
+
+	return ret;
 }
 
 static int ctrl_cmd_new_server(struct sockaddr_qrtr *from,
@@ -554,6 +568,7 @@ static int ctrl_cmd_new_lookup(struct sockaddr_qrtr *from,
 	filter.service = service;
 	filter.instance = instance;
 
+	rcu_read_lock();
 	radix_tree_for_each_slot(node_slot, &nodes, &node_iter, 0) {
 		node = radix_tree_deref_slot(node_slot);
 
@@ -568,6 +583,7 @@ static int ctrl_cmd_new_lookup(struct sockaddr_qrtr *from,
 			lookup_notify(from, srv, true);
 		}
 	}
+	rcu_read_unlock();
 
 	/* Empty notification, to indicate end of listing */
 	lookup_notify(from, NULL, true);


--- a/net/qrtr/ns.c                /* generated */
+++ b/net/qrtr/ns.c
@@ -554,6 +554,7 @@ static int ctrl_cmd_new_lookup(struct sockaddr_qrtr *from,
 	filter.service = service;
 	filter.instance = instance;

+	rcu_read_lock(); // Add RCU read lock
 	radix_tree_for_each_slot(node_slot, &nodes, &node_iter, 0) {
 		node = radix_tree_deref_slot(node_slot);

@@ -568,6 +569,7 @@ static int ctrl_cmd_new_lookup(struct sockaddr_qrtr *from,
 			lookup_notify(from, srv, true);
 		}
 	}
+	rcu_read_unlock(); // Add RCU read unlock

 	/* Empty notification, to indicate end of listing */
 	lookup_notify(from, NULL, true);
\end{minted}
\end{code}

\paragraph{Example D: \texttt{qrtr.c} port validation.}
In Listing~\ref{lst:qrtr_misfix}\footnote{Example D in Syzkaller: \url{https://syzkaller.appspot.com/bug?id=ca2299cf11b3e3d3d0f44ac479410a14eecbd326}}, the developer replaces \texttt{idr\_alloc()} with \texttt{idr\_alloc\_u32()} to avoid casting the (possibly large) \texttt{u32} port number to \texttt{int}.
By contrast, \cresearcher{} adds defensive checks that reject ports with \texttt{port < 0}, both in \texttt{qrtr\_port\_assign} and \texttt{\_\_qrtr\_bind}.
This patch resolves the crash, but rejects certain port numbers rather than handling them, so is not equivalent to the developer patch and is \emph{inaccurate}.
But the incoming value originates from \texttt{\_\_u32 sq\_port}, and special constants like \texttt{QRTR\_PORT\_CTRL} (defined as \texttt{0xfffffffeu}) are valid and widely used in the subsystem.


\begin{code}
\captionof{listing}{Developer and inaccurate patches for \texttt{qrtr.c}.}
\label{lst:qrtr_misfix}
\begin{minted}[linenos,breaklines,xleftmargin=1.5em,numbersep=5pt]{diff}
--- a/net/qrtr/qrtr.c                         /* developer */
+++ b/net/qrtr/qrtr.c
@@ -692,23 +692,25 @@ static void qrtr_port_remove(struct qrtr_sock *ipc)
  */
 static int qrtr_port_assign(struct qrtr_sock *ipc, int *port)
 {
+	u32 min_port;
 	int rc;

 	mutex_lock(&qrtr_port_lock);
 	if (!*port) {
-		rc = idr_alloc(&qrtr_ports, ipc,
-			       QRTR_MIN_EPH_SOCKET, QRTR_MAX_EPH_SOCKET + 1,
-			       GFP_ATOMIC);
-		if (rc >= 0)
-			*port = rc;
+		min_port = QRTR_MIN_EPH_SOCKET;
+		rc = idr_alloc_u32(&qrtr_ports, ipc, &min_port, QRTR_MAX_EPH_SOCKET, GFP_ATOMIC);
+		if (!rc)
+			*port = min_port;
 	} else if (*port < QRTR_MIN_EPH_SOCKET && !capable(CAP_NET_ADMIN)) {
 		rc = -EACCES;
 	} else if (*port == QRTR_PORT_CTRL) {
-		rc = idr_alloc(&qrtr_ports, ipc, 0, 1, GFP_ATOMIC);
+		min_port = 0;
+		rc = idr_alloc_u32(&qrtr_ports, ipc, &min_port, 0, GFP_ATOMIC);
 	} else {
-		rc = idr_alloc(&qrtr_ports, ipc, *port, *port + 1, GFP_ATOMIC);
-		if (rc >= 0)
-			*port = rc;
+		min_port = *port;
+		rc = idr_alloc_u32(&qrtr_ports, ipc, &min_port, *port, GFP_ATOMIC);
+		if (!rc)
+			*port = min_port;
 	}
 	mutex_unlock(&qrtr_port_lock);

--- a/net/qrtr/qrtr.c                /* generated */
+++ b/net/qrtr/qrtr.c
@@ -695,7 +695,9 @@ static int qrtr_port_assign(struct qrtr_sock *ipc, int *port)
 	int rc;

 	mutex_lock(&qrtr_port_lock);
-	if (!*port) {
+	if (*port < 0) {
+		rc = -EINVAL;
+	} else if (!*port) {
 		rc = idr_alloc(&qrtr_ports, ipc,
 			       QRTR_MIN_EPH_SOCKET, QRTR_MAX_EPH_SOCKET + 1,
 			       GFP_ATOMIC);
@@ -759,6 +761,9 @@ static int __qrtr_bind(struct socket *sock,
 		return 0;

 	port = addr->sq_port;
+	if (port < 0)
+		return -EINVAL;
+
 	rc = qrtr_port_assign(ipc, &port);
 	if (rc)
 		return rc;
\end{minted}
\end{code}